\chapter{Extracci� d'Informaci�}
\label{ie}

En aquest cap�tol tractarem

En el nostre context, anomenarem \textit{wrapper} a un tro� de codi que podem utilitzar per extreure una pe�a d'informaci� concreta d'un document.

Podem imaginar-ho com filtre que nom�s ens deixa veure una part del document que ens interessa.

\begin{figure}[h!]
\begin{center}
\includegraphics[scale=0.4]{figures/wrapper_sample.pdf}
\end{center}
\end{figure}

\section{\textit{Wrappers} a m�}

\section{Inducci� de \textit{wrappers}}

\subsection{Generaci� autom�tica de regles}
\subsubsection{Ruta d'un element HTML}
\subsubsection{Expressi� regular}

\subsection{Avaluaci� dels \textit{wrappers}}
Una vegada hem generat el conjunt dels \textit{wrappers} possibles per a un conjunt de documents, cal que avaluem quins d'ells funcionen millor. Utilitzem un sistema de vots positius i negatius i en calculem la mitjana amb la seg�ent f�rmula:

\begin{equation*}
    score = \frac{vots\; positius}{vots\; totals}
\end{equation*}

\subsection{Reaprenentatge}
El sistema est� dissenyat per tal que, quan hi ha una davallada en el nombre de refer�ncies extretes correctament, provi de reaprendre els \textit{wrappers} autom�ticament a partir dels exemples que t� emmagatzemats d'execucions passades.
